\documentclass[11pt]{article}

\usepackage[margin=1.05in]{geometry}
\usepackage{amsmath,amssymb,amsthm}
\usepackage{booktabs}
\usepackage{array}
\usepackage{microtype}
\usepackage{xcolor}
\usepackage[colorlinks=true,linkcolor=blue!60!black,citecolor=blue!60!black,urlcolor=blue!60!black]{hyperref}

\newcommand{\PT}{\mathcal{PT}}
\newcommand{\ee}{\mathrm{e}}
\newcommand{\ii}{\mathrm{i}}
\newcommand{\norm}[1]{\left\lVert #1 \right\rVert}

\newtheorem{theorem}{Theorem}
\newtheorem{proposition}{Proposition}
\newtheorem{corollary}{Corollary}
\newtheorem{lemma}{Lemma}
\theoremstyle{remark}
\newtheorem{remark}{Remark}

\title{The conformal point is a phase boundary:\\
unfolding directions, the sign map of the cascade divisor,\\
and a two-channel photonic readout}
\author{Jonathan Torkelson\\[2pt]
\small Embeddetech, Inc., Tulsa, Oklahoma, USA \quad
\texttt{jtorkelson@embeddetech.com}}
\date{July 12, 2026}

\begin{document}
\maketitle

\begin{abstract}
\noindent The equal-frequency Pais--Uhlenbeck point---the per-mode
kinematics of conformal gravity---is an exceptional point (EP), and
this paper treats it as what an exceptional point is: a phase
boundary. We classify every deformation of the theory by two data:
the \emph{kind} of unfolding direction (family-parameter detunings
split the degenerate pair linearly; generic matrix perturbations
split by Puiseux $\sqrt\epsilon$) and its \emph{phase}, which
decides oscillatory versus growing dynamics. On this map the
oscillatory set is the punctured real axis ($\Delta\neq0$) of the family directions (while
$\omega_2^2$ stays positive) and, the ray $\varphi=0$ (the
single-entry-model reference direction, resolved to within $O(\epsilon)$) of the
generic ones---measure zero (in the complex plane) either
way (the family real axis by the analytic law; the generic ray by
the exact Puiseux eigenvalues, from which growth with probability
one follows for any absolutely continuous phase law (unstructured, traceless/phase-only regulators; fully generic noise still gives probability $\to1$, at rate $1-O(\sqrt\epsilon)$ for phase densities bounded near the band)---corroborated by
random-phase perturbations landing on the
growing side in $500/500$ draws at $\epsilon=10^{-8}$),
so any \emph{unstructured} regulator of the pure theory buys the
growing side. Real-symplectic (structure-preserving) regulators are the
sole exception---the characteristic polynomial stays even and real, so
its discriminant sign splits perturbation directions into two
\emph{open} half-spaces, an oscillatory band and a growing one, and
which side is bought then has probability strictly between $0$ and $1$
(developed below). The known deformations take their places: an Einstein term
unfolds on the real axis (Stelle, oscillatory, dictionary cost
$2/\delta$); the Donoghue--Menezes loop width tilts the approach ray
by a fixed angle without imposing a detuning floor on the approach to
the exceptional line; and the
Ilhan--Kovner ``ghost-turned tachyon'' blowup of conformal gravity's
$k=0$ sector is the corner case $\omega\to0$, where the family
direction itself turns Puiseux---their finite-time singularity sits on
the corner's imaginary-split unfolding (the $\varphi=\pi$ direction,
where the detuned $\omega_2^2<0$); their blowup is reproduced here
with no free parameters, while the map's frozen-mode rate
$(\sqrt3/2)\gamma_b$ is the rate their equations imply once their
effective detuning is ingested.
The boundary is then shown to be readable at the level of the
theory's stochastic (beable) representation: for cascade
(triangular) processes we prove the divisor of the process is
triangular with a per-bond sign law; that within Jordan cascades
the pure conditional-mass-creation mechanism of the companion paper
is strictly a property of the minimal block---for Jordan blocks of
size $k\ge3$ the divisor generator has entries of \emph{both}
signs, in an exactly alternating pattern given by the Rayleigh sums
of the zeros of the Bessel function $J_0$. At finite steps a
dominant-pole argument---asymptotic, with numerically computed
window edges, below the strength of the theorems---locates
the divisor's negative windows on Fourier--Bessel
bands with edges at $(j_{0,n}/j_{0,1})^2$. Real detuning brings
negativity back at the per-bond first-return length
$\pi/|\delta_i|$; for splits in the closed upper half-plane the
recurrence is paced by the \emph{real part} of the split, while
the negativity-free directions comprise the $+\ii$ ray, the EP,
and a computed silent cone centered on the $-\ii$ ray whose sharp
thresholds reach to within ${\approx}15.3^\circ$ of the real
axis, from below (silent split phases ${\approx}195.3^\circ$--$344.7^\circ$; thresholds
located to $10^{-3}$ degrees, cone interior sampled, in the text): the
negativity channel is a spectrometer for the magnitude of the real
part of the split, with a known silent sector. The divisor's second
diagnostic---the column-sum (conditional-mass-creation) channel,
stated relative to the minimal contraction normalization defined
in the text---fires at every sampled direction (a $5^\circ$ scan
of the full circle finds no silent phase) and its onset, at known
$|\delta|$, resolves the growing from the decaying imaginary ray
(relative growth of the fed level, under the symmetric split
convention fixed in the text)---a sign the negativity channel
cannot see; and the $\omega\to0$ corner, being a size-4 Jordan
block rather than a bond, is negativity-bearing by the $k\ge3$
theorem at the corner itself and, computationally, at every
sampled unfolding direction. The two channels together leave no
sampled direction silent, bond by bond. A two-device photonic
extension of our measurement proposal (a coupler \emph{at} its EP,
and a dissipatively-coupled unidirectional pair) would measure the
mechanism fingerprint and the direction spectroscopy on
capabilities demonstrated in the platform family, one coupling
transfer priced into the device error budget. All results are tree-level mode kinematics, stated for
the minimal cell and its finite stackings.
\end{abstract}

\newpage
\begin{quote}
\small\noindent\textbf{Note on provenance.} This manuscript is an
experiment in AI-conducted research as much as a physics paper. The
project began when the author asked the large language model Claude
Fable~5 (Anthropic): \emph{``Is it not possible that the antilinear
PT/time-reversal symmetry element could maybe predict indivisible
non-Markovian processes?''} Everything that followed was performed by
the model: the research program, the theorems and proofs, the numerical
experiments, the literature search, and this text. The model was
repeatedly prompted to attack its own results adversarially and to
search for prior work. The author, a professional engineer rather than
a physicist, directed and reviewed the sessions but makes no claim to
comprehension of the findings, and is circulating the manuscript for
independent review by professional physicists. Accordingly, nothing
here should be taken on trust: every theorem is stated with a full
proof, and every numerical claim is reproducible from the public
scripts, with one disclosed exception (the Merlin row of
Sec.~\ref{sec:dictionary}, which is analytic---its arithmetic stated in
the revised companion; see Code availability); the originating
conversation,
project briefing, and full research notebook are preserved in the
program's hub repository,
\url{https://github.com/embeddetech/ghosts-program}.
\end{quote}

%=====================================================================
\section{Introduction}
\label{sec:intro}

This is the sixth paper of a program
that reads quantum dynamics in
Barandes' currency of indivisible stochastic processes on
configurations~\cite{Barandes} and asks what becomes of ghost
sectors there~\cite{indivisible,instantaneous,conditioned,choosing}. The
conformal-gravity paper~\cite{choosing} established the arena for the present one: the
equal-frequency Pais--Uhlenbeck (PU) mode pair---the kinematics that
conformal gravity carries at every spatial momentum---is a Jordan
block, an exceptional point (EP) of the mode dynamics, and at that
point no closed stochastic representation of any kind exists, while
the open, survival-conditioned representation is a one-way cascade
whose indivisibility enters through a new mechanism (\emph{conditional
mass creation}: the two-point divisor stays entrywise nonnegative but
its column sum crosses one, a diagnostic stated relative to the
minimal contraction normalization made explicit in
Sec.~\ref{sec:divisor})~\cite{choosing}.

That paper treated the point itself. This paper treats its
neighborhood, and the organizing claim is the title's: \emph{the
conformal point is a phase boundary}, and everything in its vicinity
is classified by how a deformation unfolds the Jordan block. The
claim has three layers.

\emph{Layer 1: the direction map} (Sec.~\ref{sec:map}). Deformations
come in two kinds. Physical parameter changes---an Einstein--Hilbert
term, a loop-induced width---move within the PU \emph{family},
detuning $\omega_2^2=\omega^2+\Delta$ and splitting the degenerate
frequencies \emph{linearly} in $|\Delta|$ (the non-generic unfolding
responsible for the $2/\delta$ dictionary cost computed
in~\cite{choosing}). Generic \emph{unstructured} perturbations---floating-point
noise and random regulators that touch
the matrix rather than the family (a structured real-symplectic lattice
is the exception treated below, with an open oscillatory band)---split by Puiseux
$\sqrt\epsilon$~\cite{Kato}. In both kinds the \emph{phase} of the
perturbation decides the dynamics, with a kind-dependent
oscillatory set (within an $O(\epsilon)$ neighborhood of the positive-real ray for traceless perturbations;
the punctured real axis ($\Delta\neq0$), while $\omega_2^2>0$, for the family)
and growth everywhere else (for net-damping fully generic noise an $O(\sqrt\epsilon)$-wide decaying band intervenes, so growth there holds only with probability $1-O(\sqrt\epsilon)$). The oscillatory set is measure zero in
the generic directions, and the $\omega\to0$ corner---conformal
gravity's homogeneous sector---is where even the family direction
turns Puiseux: the Ilhan--Kovner instability~\cite{IK2014} is
recovered there, quantitatively, as the corner's imaginary-split
unfolding.

\emph{Layer 2: the divisor image} (Sec.~\ref{sec:divisor}). The
program's currency sees the same boundary. For cascade processes we
prove a per-bond sign law for the two-point divisor and show that
the sign structure of a deep Jordan cascade's divisor is universal
Bessel-zero arithmetic: the divisor generator's entries alternate in
sign with coefficients given exactly by the Rayleigh sums
$R_p=\sum_n j_{0,n}^{-2p}$, and finite-step negative windows
converge to Fourier--Bessel bands (a dominant-pole asymptotic,
below theorem strength). Two corollaries matter
physically: the mass-creation-pure mechanism of~\cite{choosing}
is confined to the minimal ($k=2$) Jordan cell---no taller Jordan
cascade is negativity-free---and a real-detuned cascade's
divisor develops negativity at the first-return length
$\pi/|\delta_i|$ per bond, the recurrence for upper-half-plane
splits being paced by the real part of the split (a cone about
the $-\ii$ ray is silent for negativity; Sec.~\ref{sec:divisor}).
The negativity channel alone is therefore a spectrometer for the
magnitude of the \emph{real} part of the split; its silent set---the
$+\ii$ ray, the EP, and the cone about the $-\ii$ ray---is
read out at every sampled phase by the divisor's second, convention-relative
channel (the column sum, silent at no sampled phase), whose onset resolves the
sign negativity cannot see, and the
$\omega\to0$ corner---a size-4 block---is negativity-bearing by
the $k\ge3$ theorem before any unfolding is chosen
(Sec.~\ref{sec:divisor}): a two-channel direction readout.

\emph{Layer 3: the readout} (Sec.~\ref{sec:photonic}). Both the
mechanism fingerprint and the direction spectroscopy are measurable
with classical light on the femtosecond-laser-written waveguide
platform of Klauck et al.~\cite{Klauck2019}, via a two-device
extension of our standing proposal: a coupler operated exactly at
its EP, and a dissipatively-coupled pair at the unidirectional point
that realizes the PU cascade cell in the guide basis.

Throughout, the honest scope is the one inherited
from~\cite{choosing}: these are tree-level statements about mode
kinematics, proved for the minimal cell and finite stackings, with
the interacting field theory entering only through the
deformation dictionary (Sec.~\ref{sec:dictionary}).

%=====================================================================
\section{The direction map}
\label{sec:map}

Write the mode pair in its normal form. At the conformal point the
one-particle matrix of a mode is similar to the $2\times2$ Jordan
cell $J=\omega\mathbf 1+N$, $N_{12}=1$. A deformation acts either
\emph{within the PU family}---the physical parameters
$(\omega_1^2,\omega_2^2)$ move, $\omega_2^2=\omega^2+\Delta$ with
$\Delta=|\Delta|\ee^{\ii\varphi}$---or \emph{generically}, as a
matrix perturbation, modeled by the block entry
$H=\bigl(\begin{smallmatrix}\omega&1\\ \epsilon\ee^{\ii\varphi}&
\omega\end{smallmatrix}\bigr)$.

\paragraph{Family direction: linear split.} The eigenfrequencies of
the detuned pair are $\omega$ and
$\sqrt{\omega^2+\Delta}\simeq\omega+\Delta/2\omega$: the split is
linear in $|\Delta|$ (the family direction is a non-generic
unfolding; this is why the dictionary cost of~\cite{choosing}
scales as $2/\delta$ rather than $2/\sqrt\delta$, re-verified below).
The growth rate is
\begin{equation}
\lambda_{\rm fam}(\Delta)=\bigl|\mathrm{Im}\sqrt{\omega^2+\Delta}
\bigr|\;\simeq\;\frac{|\Delta|\,|\sin\varphi|}{2\omega},
\label{eq:famlaw}
\end{equation}
verified on the $4\times4$ PU companion matrix at
$|\Delta|=10^{-3},10^{-2},10^{-1}$, $\varphi\in\{0,\pi/6,\pi/2,
5\pi/6,\pi\}$---the lower half-plane follows by conjugation
symmetry, $\lambda(\bar\Delta)=\lambda(\Delta)$ (worst relative
deviation $4.6\times10^{-2}$ over the nonzero-growth phases,
dominated by
the $|\Delta|=10^{-1}$ curvature; Eq.~\eqref{eq:famlaw} is the
$|\Delta|\to0$ law). Two features deserve emphasis. First, the
oscillatory set of the family direction is the punctured \emph{real
axis} of $\Delta$ ($\Delta\neq0$, both signs), so long as $\omega_2^2>0$: at
$\varphi=0$ and $\varphi=\pi$ the computed growth is
$\lesssim10^{-13}$. Second, on the real axis the EP approach carries
the known cost: the eigenvector-matrix condition number obeys
$\mathrm{cond}(S)\,\delta/2=1.008$ and $1.001$ at
$\delta=5\times10^{-3},\,5\times10^{-4}$---the $2.00/\delta$ law
of~\cite{choosing}, recovered here as the boundary-line toll.

\paragraph{Generic direction: Puiseux split.} The perturbed block
has eigenvalues $\omega\pm\sqrt{\epsilon}\,\ee^{\ii\varphi/2}$, so
\begin{equation}
\lambda_{\rm gen}(\epsilon)=\sqrt{\epsilon}\,\bigl|\sin(\varphi/2)
\bigr|,
\label{eq:genlaw}
\end{equation}
verified to $1.2\times10^{-13}$ across the grid, with the fitted
Puiseux exponent $0.5000$. Here the oscillatory set is the single
ray $\varphi=0$ of this single-entry model (a fully generic matrix
perturbation adds a common $O(\epsilon)$ growth $\tfrac12\epsilon\,
\mathrm{Im}\,\mathrm{tr}\,P$, absent here where $\mathrm{tr}\,P=0$; the
oscillatory locus is instead displaced from this ray by only $O(\epsilon)$
through the sub-leading discriminant correction (and is empty once $\mathrm{Im}\,\mathrm{tr}\,P\neq0$, its role taken by the
growing/decaying split below) and, being contained in the zero set of a single real-analytic
condition (the imaginary part of the discriminant, not identically
vanishing in the phase), is measure zero a fortiori): a random-phase perturbation at $\epsilon=10^{-8}$
landed on the growing side in $500$ of $500$ draws, with median
rate $\sim\sqrt\epsilon$. \emph{The pure theory is the boundary}
is thus an analytic statement---growth with probability one for
any absolutely continuous phase law, from the exact Puiseux
eigenvalues---corroborated by sampling: an \emph{unstructured} regulator
that perturbs the block with a continuously-distributed phase
(floating-point roundoff, modeled as an effectively random phase over inputs,
being the representative case) buys the growing
side with probability one, at the Puiseux-enhanced rate $\sqrt\epsilon$;
a fully generic (non-traceless) perturbation instead opens, when the common
term damps ($\mathrm{Im}\,\mathrm{tr}\,P<0$), a decaying band of
angular width $O(\sqrt\epsilon)$ through the common
$\tfrac12\epsilon\,\mathrm{Im}\,\mathrm{tr}\,P$ term, so there the growing
side has probability $\to1$ as $\epsilon\to0$ (at rate $1-O(\sqrt\epsilon)$ for
phase densities bounded near the band).
Structured regulators are a different matter: a real-symplectic lattice
discretization sits at a defective $1{:}1$ resonance, whose Hamiltonian
unfoldings keep an open set of directions oscillatory (a property of the
real-symplectic unfolding, not of a Krein signature---undefined at the
defective point~\cite{KreinTheory}), so there, for any direction law with full
support, the growing side has probability
strictly between zero and one (an open band of oscillatory directions
survives); and a deterministic truncation lands at one definite
phase, for which ``probability one'' is not the right
description.

\paragraph{The $\omega\to0$ corner.} Conformal gravity's
translationally-invariant ($k=0$) sector sits at the doubly
degenerate corner $\omega_1=\omega_2=0$, with secular solution set
$\{1,t,t^2,t^3\}$. There the family direction itself becomes
Puiseux: the induced frequencies are $\pm\sqrt{\Delta}$, so
$\lambda_{\rm corner}=\sqrt{|\Delta|}\,|\sin(\varphi/2)|$ (verified
to $6.9\times10^{-16}$). The two phases in play are related by the
square root: on the finite-frequency line the split is
$\delta\simeq\Delta/2\omega$ and shares the phase of $\Delta$,
while at the corner $\delta=\sqrt\Delta$ \emph{halves} it---which
is how a $\varphi=\pi$ family deformation, real in $\Delta$,
arrives as a purely imaginary frequency split (the $+\ii$ ray of
Sec.~\ref{sec:divisor}'s split-phase classification, on the
principal branch; the corner block itself sits outside the
bond-level classification, as Sec.~\ref{sec:divisor} shows). This corner
is where the Ilhan--Kovner result lives. In~\cite{IK2014} the nonlinear
self-interaction of conformal gravity's $k=0$ sector generates, for
a background amplitude $\gamma_b$, an effective potential whose
sign relative to the ghost kinetic term is anti-restoring: in our
variables, $\Delta_{\rm eff}=-\tfrac34\gamma_b^2$, the
$\varphi=\pi$ ray. With their $\Delta_{\rm eff}$ ingested, the map
returns growth at
$\sqrt{3/4}\,\gamma_b\cdot|\sin(\pi/2)|=(\sqrt3/2)\gamma_b$---algebra
once $\Delta_{\rm eff}$ is given, not an independent
prediction---and the frozen-mode eigenvalue check confirms rates
$0.433013/0.866025/1.732051$ at $\gamma_b=0.5/1/2$, matching
$(\sqrt3/2)\gamma_b$ to the six digits shown. The parameter-free
corroboration is separate: direct integration of their
\emph{nonlinear} equations of motion (public script, block~[3b])
reproduces their exact solution, Eq.~(58) of~\cite{IK2014}, to
relative $5\times10^{-12}$ and the singularity time
$t_{\rm sing}=2^{3/2}/\gamma_0$ ($\gamma_0$ their initial amplitude) to
relative $\sim10^{-6}$ (five--six significant digits).
Their ``ghost-turned tachyon'' is the
imaginary-split unfolding of the conformal corner (a $\varphi=\pi$
family direction whose frequency split is imaginary because
$\omega_2^2<0$ there); the reproduction, not the rate, is what
carries ``no free parameters.''

Three scope guards, which we owe to a careful reading
of~\cite{IK2014,IK2013}. Their exceptional point is the
$\omega=0$ corner, not the finite-frequency line: the extension of
the classification to finite $k$ is ours, not theirs. Their
detuning is \emph{self-induced} and amplitude-dependent
($\Delta_{\rm eff}\propto\gamma_b^2$), so the unfolding statement
is about the interacting background, never about the linear
spectrum. And the companion paper~\cite{IK2013} exhibits an
interacting particle--ghost system that remains stable---in our
language, an interaction that unfolds along the real
direction---so ``interactions destabilize ghosts'' is not a
universal; \emph{which way the interaction unfolds the block} is
the invariant content, and their stable example is a test case the
classification passes rather than an exception it suffers.

%=====================================================================
\section{The divisor image}
\label{sec:divisor}

The companion papers measure dynamics in the beable currency: a
mode-pair evolution $U(t)$ is read through
$\Gamma(t)=|U(t)/\sigma_{\max}(t)|^2$ (entrywise moduli of the
contraction-normalized evolution), and divisibility is interrogated
through the two-point divisor $M=\Gamma(s+\tau)\,\Gamma(s)^{-1}$.
The indivisible-process paper~\cite{indivisible} found the broken
dimer fails sub-stochasticity by an entry going \emph{negative}; the
conformal-gravity paper~\cite{choosing} found the PU cell fails by the
column sum crossing one with nonnegativity intact (\emph{conditional
mass creation}). This section proves
where each mechanism lives.

\paragraph{Conventions.} Three conventions are inherited
from~\cite{indivisible,choosing} and restated here, because the
results below split cleanly by whether they depend on them. The
\emph{beable basis} is the configuration basis in which the cascade
Hamiltonian is triangular; processes are read entrywise,
$\Gamma=|{\cdot}|^2$; and $\sigma_{\max}(t)$ denotes the largest
singular value of $U(t)$, so that $C(t)=U(t)/\sigma_{\max}(t)$ is
the \emph{minimal contraction normalization}: a unitary
(conservative) dilation exists exactly when the normalized evolution
is a contraction, so $\sigma_{\max}$ is the pointwise-least
normalization admitting a conservative completion---the unique
choice that discards no survival probability beyond what the
dynamics itself loses ($\norm{C(t)}=1$: at every time the
least-decaying state is not conditioned at all). The two failure
diagnostics then differ in kind. Divisor entry \emph{signs} are
invariant under every positive scalar renormalization $c(t)$ of the
normalization---negativity (M1) is convention-free, as
Proposition~\ref{prop:signlaw}(ii) records---while divisor
\emph{column sums} scale by $c^2(s)/c^2(s+\tau)$: the
column-sum-crossing diagnostic (conditional mass creation, M2) is
stated relative to the $\sigma_{\max}$ convention, and its
onset---indeed its occurrence---moves under renormalization. The
relativity is quantifiable, and quantifying it is what justifies
the convention: for the unit cell with coupling $\kappa$ the raw
$|U|^2$ divisor's column sum over the step from $s$ to $T=s+\tau$
is exactly $1+\kappa^2(T^2-s^2)$, so
crossing-freedom at the pair $(s,T)$ forces
$c^2(T)/c^2(s)\ge1+\kappa^2(T^2-s^2)$; telescoping this over a
refined partition of $[0,t]$ turns the product into the Gaussian, so
a renormalization avoids the crossing at \emph{all} step pairs only
if $c^2(t)\ge c^2(0)\,\ee^{\kappa^2t^2}$, and that Gaussian envelope
does avoid it (the exact column-sum law and the envelope's
sufficiency verified in \texttt{scripts/p6\_divisor\_gate.py},
block~[4]; the necessity is the telescoping just given). Divisibility
can thus be purchased---within positive scalar renormalizations
$c(t)$---but only by
conditioning away survival probability at a Gaussian rate against
polynomially growing dynamics---a conditioning overhead
$c^2/\sigma_{\max}^2$ that diverges; under the minimal convention,
in which conditioning means survival and nothing else, the crossing
is unavoidable. Convention-free content below: the triangularity of
the process and every divisor sign statement
(Proposition~\ref{prop:signlaw}, Theorems~\ref{thm:k2}
and~\ref{thm:rayleigh}). Convention-relative: every mass-creation
onset ($t_c$, $z_c$).

\paragraph{Notation.} From here on, $\delta$ (with per-bond
subscripts $\delta_i$, and as the device detuning of
Sec.~\ref{sec:photonic}) always denotes a \emph{frequency} split,
as in Sec.~\ref{sec:map}; the symbol $\Delta$ is reserved for
Sec.~\ref{sec:map}'s family detuning, which lives in $\omega^2$
($\omega_2^2=\omega^2+\Delta$). On the finite-frequency line
$\delta\simeq\Delta/2\omega$; at the corner $\delta=\sqrt\Delta$
halves the phase.

\begin{proposition}[cascade divisor; per-bond sign law]
\label{prop:signlaw}
Let $H$ be upper triangular on $d$ configurations (a cascade), and
let $V(t)=|U(t)|^2$ entrywise, $U(t)=\ee^{-\ii Ht}$. Then for all
$0<s<s+\tau$: (i) the divisor $M=V(s+\tau)V(s)^{-1}$ is upper
triangular with strictly positive diagonal
$M_{ii}=V_{ii}(s{+}\tau)/V_{ii}(s)$; and (ii) for each bond $i$,
\[
\operatorname{sign} M_{i,i+1}
=\operatorname{sign}\bigl[r_i(s+\tau)-r_i(s)\bigr],
\qquad
r_i(t)=\frac{|U_{i,i+1}(t)|^2}{|U_{ii}(t)|^2},
\]
independently of the $\sigma_{\max}$ normalization.
\end{proposition}

\begin{proof}
$U(t)$ is upper triangular, hence so are $V$, $V^{-1}$, and $M$;
the diagonal of $M$ is the ratio of positive diagonals. For (ii),
the inverse of an upper-triangular matrix has
$(V^{-1})_{i,i+1}=-V_{i,i+1}/(V_{ii}V_{i+1,i+1})$, so
\[
M_{i,i+1}
=\frac{V_{i,i+1}(s{+}\tau)}{V_{i+1,i+1}(s)}
-\frac{V_{ii}(s{+}\tau)\,V_{i,i+1}(s)}{V_{ii}(s)V_{i+1,i+1}(s)}
=\frac{V_{ii}(s{+}\tau)}{V_{i+1,i+1}(s)}
\bigl[r_i(s{+}\tau)-r_i(s)\bigr],
\]
a positive multiple of the ratio increment. The $\sigma_{\max}$
normalization multiplies $M$ by a positive scalar.
\end{proof}

Adversarial check: $1408/1408$ sign agreements over random complex
triangular families at $d=3,4,5$, with the lower-triangular
residual exactly zero.

For the detuned $2\times2$ cascade
$H=\bigl(\begin{smallmatrix}\omega_1&c\\0&\omega_2
\end{smallmatrix}\bigr)$ the ratio is explicit:
$r(t)=|c|^2\bigl(2-2\cos\delta t\bigr)/\delta^2$ for real split
$\delta=\omega_1-\omega_2$ (oscillates: negativity \emph{recurs},
first onset $s^*=\pi/|\delta|$; computed onset ratios
$0.9969$--$0.9995$ over $\delta=0.4$--$2.0$, at a $0.02$ divisor
step with a $0.01$ base-point grid---the detected onset is the
first grid point past $\pi/|\delta|-\tau/2$);
$r(t)=|c|^2\bigl(1-\ee^{(g_1-g_2)t}\bigr)^2/(g_1-g_2)^2$ for a
decay split (monotone: negativity-free at all times, machine
epsilon); and $r(t)=|c|^2t^2$ at the EP (monotone: the
mass-creation-pure cell of~\cite{choosing}, onset re-verified at
$0.48$ at a $0.05$ pair spacing with a $0.01$ base grid; the
zero-step limit is $0.486$, of which the companion's $\approx0.50$
was its coarser $0.05$ base grid rounding up---finite-pair
conventions can shift the detected onset in either direction). The general split $\delta=a+\ii b$ interpolates:
$r(t)=|c|^2|\ee^{\ii\delta t}-1|^2/|\delta|^2$, whose derivative
is $\propto\ee^{-bt}\bigl[b\cos at+a\sin at-b\,\ee^{-bt}\bigr]$.
For $b\ge0$ (splits in the closed upper half-plane) this changes
sign whenever $a\neq0$, so there the negativity-free directions
are exactly the $+\ii$ ray and the EP, and \emph{the negativity
channel reads out the magnitude of the real part of the split}:
real-direction unfolding brings
negativity back with first-return time $\pi/|\delta|$, diverging as
the EP is approached, and a mixed split recurs on the pace of
$|\mathrm{Re}\,\delta|$ (its sign is a residual ambiguity of the
negativity channel). For $b<0$ the term $-b\,\ee^{-bt}
=|b|\ee^{|b|t}$ outgrows the oscillation amplitude
$\sqrt{a^2+b^2}$ at $t^*=\ln\bigl(\sqrt{a^2+b^2}/|b|\bigr)/|b|$,
and when $t^*$ precedes the first trough the ratio is, heuristically,
monotone despite $a\neq0$ (the thresholds below come from numerical
bisection, not from this sufficient condition): the negativity-free set acquires a
\emph{silent cone} (interior sampled) centered on the $-\ii$ ray, spanning split
phases $195.345^\circ$--$344.655^\circ$ (sharp, scale-invariant
thresholds within $15.345^\circ$ of the real axis (from below); bisection
to $10^{-3}$ degrees; at the thresholds themselves the derivative
is tangent to zero, so they are silent). The wedges between the
real axis and the thresholds are negativity-bearing but
\emph{transient}: finitely many windows before the growth term
silences the bond---a single window over the few degrees adjacent
to the threshold, proliferating toward the real axis, where the
silencing time $t^*$ diverges (computed, at the sampled wedge
phases: $1$ window at $190^\circ$ and $193^\circ$, $2$ at
$188^\circ$---completed totals, scale-invariant---and $4$ within
$t\le40$ at $|\delta|=0.6$ at $181^\circ$ and $185^\circ$). A wedge direction is thus
readable only in a finite early epoch, in contrast with the
indefinitely recurrent upper half-plane.
Computed over the full circle at
$|\delta|=0.6$, $t$ on the uniform grid $0.05{:}0.02{:}40$:
monotone ($0$ decreasing intervals) at $90^\circ$ and at the
sampled phases
$210^\circ$--$330^\circ$ (interior to the computed silent cone, whose
connectedness we read from this sampled grid rather than a dense scan),
versus several hundred to ${\sim}10^3$ decreasing intervals at the
sampled phases $0^\circ,30^\circ,60^\circ,120^\circ,150^\circ,
180^\circ$ ($919$ at $30^\circ/150^\circ$, $953$ at
$0^\circ/180^\circ$; the near-threshold count at
$60^\circ/120^\circ$ is tolerance-sensitive in the exponentially
flattened tail, its implementation-robust strict count being
$1050$). The negativity channel is
thus a spectrometer for the magnitude of the real part of the split,
with a known
silent sector. In a
real-detuned bidiagonal chain the readout is per
bond: each bond $i$ returns at $\pi/|\delta_i|$ (computed ratios
$0.9982$--$0.9995$, $d=4$).

\paragraph{The second channel, and the corner.} The silent
directions of the negativity channel are not silent for the
divisor as a whole. The column-sum diagnostic---conditional mass
creation, convention-relative in the sense of the conventions
paragraph---fires at every \emph{sampled} split phase: computed
onsets
over the full circle at $|\delta|=0.6$ (a $5^\circ$ grid, $72$
phases; coupling $1$, pair step
$0.05$, base grid $0.01$; \texttt{scripts/p6\_m2\_channel.py}) range from
$0.30$ to $0.68$ with no silent phase found, and on the imaginary
axis
the onset is \emph{asymmetric}---$0.68$ on the growing ($+\ii$)
ray against $0.30$ on the decaying ($-\ii$) ray ($0.58$ vs $0.38$
at $|\delta|=0.3$; $0.69$ vs $0.30$ at pair step $0.02$)---so the
second channel, at known $|\delta|$, resolves the sign of the
\emph{imaginary} part (the $\pm\ii$ ray); both channels are even under $\mathrm{Re}\,\delta\to-\mathrm{Re}\,\delta$,
so the sign of the \emph{real} part stays invisible to either. With the split assigned
symmetrically, $\omega_{1,2}=\omega\pm\delta/2$ (as in the scripts;
the beable process is invariant under a common frequency shift, so
$\delta$ alone does not fix an absolute assignment), the $+\ii$ ray
is the direction in which the fed (downstream) level grows relative
to its source; ``growing/decaying ray'' is shorthand for this
relative sense. And the flagship corner instability of Sec.~\ref{sec:map}
never relied on the bond-level statement in the first place: the
$\omega\to0$ sector is a size-4 Jordan block, not a bond, so
Theorem~\ref{thm:k2} applies, and its divisor is
negativity-bearing at depth two for every sufficiently short
step---by the theorem for the corner block itself, and
computationally at every sampled unfolding direction. Computed on
the corner companion matrix (base grid $0.002$), the negativity
onset sits at the first grid point past the theorem's short-step
boundary
$s^*=\tau/(\sqrt3-1)$ ($0.070$ vs $s^*=0.0683$ at $\tau=0.05$;
$0.028$ vs $0.0273$ at $\tau=0.02$), identically for the pure
corner, the $\varphi=\pi$ (Ilhan--Kovner) unfolding, and the
$\varphi=0$ real-split one---and identically, onset $0.070$, at
all twelve sampled phases of a full unfolding circle at
$|\Delta|=0.36$---while the column-sum onsets differ by
direction ($0.800$, $0.906$, $0.728$ respectively at
$\tau=0.05$).
The instrument is thus two-channel: negativity reads the real part
of the split and is silent on the imaginary directions at the bond
level; the column sum---stated in the minimal convention---fires
at every sampled direction and resolves the imaginary sign; and
the deep-block
corner announces itself through the sign theorem before any
unfolding is chosen.

Purity, however, is a property of the \emph{minimal Jordan} cell
only:

\begin{theorem}[mass-creation purity of Jordan cascades is
$k=2$-only]
\label{thm:k2}
For the Jordan block $J_k$, $k\ge3$, at any eigenvalue (a complex
eigenvalue rescales $V$ by a positive factor, leaving divisor signs
unchanged), write
$V(t)=f(t^2N)$ with $f(x)=\sum_m x^m/(m!)^2$ and $N$ the nilpotent.
The divisor is $h(N)$ with $h(x)=f(T^2x)/f(s^2x)$, $T=s+\tau$, and
its second-superdiagonal coefficient is
\[
h_2=\tfrac14\,(T^2-s^2)(T^2-3s^2)\;<\;0
\quad\text{for }1<T^2/s^2<3 .
\]
In particular every sufficiently short step ($T<\sqrt3\,s$) has a
negative divisor entry, and the divisor generator satisfies
$L_{i,i+2}(t)=-t^3$ exactly in the Jordan basis (a fixed positive multiple
$|c_ic_{i+1}|^2$ of it in a general chain, leaving the sign unchanged): no irreducible Jordan cascade of depth $\ge3$---an
equal-diagonal, nearest-neighbor bidiagonal chain, the physical object
throughout---is negativity-free.
\end{theorem}

\begin{proof}
$|U_{i,i+m}|^2=t^{2m}/(m!)^2$ for the single block, so
$V(t)=f(t^2N)$; functions of $N$ commute, giving $h$. A general
equal-diagonal chain (bonds $c_i$) is a fixed positive-diagonal
similarity $V=D^{-1}f(t^2N)D$ of the block, multiplying each divisor
entry by a positive bond factor and leaving every sign invariant. Direct
expansion of $f(bx)/f(ax)$ to second order gives
$h_2=b^2/4-ab+3a^2/4=(b-a)(b-3a)/4$ with $a=s^2$, $b=T^2$. For the
generator, $L=\dot V V^{-1}=2tN\,(f'/f)(t^2N)$ and
$[f'/f]_1=-\tfrac12$ (or see Theorem~\ref{thm:rayleigh}), so
$L_{i,i+2}=2t\cdot(-\tfrac12 t^2)=-t^3$.
\end{proof}

Verified: closed form against direct matrix divisors to
$5.9\times10^{-13}$; the window boundary located at
$T=1.73250$ against $\sqrt3=1.73205$ (grid $5\times10^{-4}$); the
generator entries $2t$ and $-t^3$ to six digits. The conformal-gravity
paper's statement is untouched---its cell is $k=2$, where the single
superdiagonal is governed by Proposition~\ref{prop:signlaw} and
$r=t^2$ is monotone---but stacked sectors do not escape it: the
$M$-quantum sector of the equal-frequency tower carries a
\emph{single, irreducible} Jordan block of size $M{+}1$ (the
sector-tower proposition of~\cite{choosing}, the constructive
form of Smilga's infinite-order exceptional
point~\cite{Smilga2014}), so the theorem extends to every sector
$M\ge2$: the sector shift $d\Gamma(N)$ carries strictly positive weights
$\sqrt{(\nu+1)(M-\nu)}$ (occupation index $\nu=0,\dots,M{-}1$), so the beable and Jordan bases differ by one fixed,
$t$-independent positive-diagonal similarity at every $M$~\cite{choosing},
under which the entrywise divisor signs are invariant: the
tower cannot present a pure mass-creation
process---negativity at depth two is a theorem at every sector $M\ge2$, given the
companion's sector-tower structure, for every sufficiently short step.

The full sign structure is Bessel-zero arithmetic:

\begin{theorem}[Rayleigh sign law for the cascade generator]
\label{thm:rayleigh}
For the Jordan cascade of any depth, at any eigenvalue (a complex
eigenvalue shifts only the diagonal of the generator, cf.\
Theorem~\ref{thm:k2}), the divisor generator's
$(m{+}1)$-th-superdiagonal entries are, in the Jordan basis (a positive
bond factor multiplies them in a general chain, leaving signs unchanged),
\[
L_{i,\,i+1+m}(t)\;=\;2\,(-1)^m\,4^{\,m+1}R_{m+1}\,t^{2m+1},
\qquad
R_p=\sum_{n\ge1} j_{0,n}^{-2p},
\]
where $j_{0,n}$ are the positive zeros of $J_0$ and $R_p$ are
the classical Rayleigh sums (exact rationals: $R_1=\tfrac14$,
$R_2=\tfrac1{32}$, $R_3=\tfrac1{192}$,
$R_4=\tfrac{11}{12288}$, \dots)~\cite{Watson}. The signs
strictly alternate at every depth and every $t>0$.
\end{theorem}

\begin{proof}
$f(x)=\sum_m x^m/(m!)^2=I_0(2\sqrt x)$ is entire of order
$\tfrac12$ with zeros exactly at $x_n=-j_{0,n}^2/4$, so its
Hadamard factorization is $f(x)=\prod_n(1+4x/j_{0,n}^2)$ and
\[
\frac{f'}{f}(x)=\sum_n\frac{1}{x+j_{0,n}^2/4},
\]
a sum of simple poles on the negative axis with unit weights.
Expanding each term geometrically,
$[f'/f]_m=(-1)^m4^{m+1}R_{m+1}$, and
$L=2tN\,(f'/f)(t^2N)$ distributes the coefficients along the
superdiagonals. Alternation is manifest since $R_p>0$.
\end{proof}

The exact rational coefficients were confirmed against direct sums
over $2\times10^5$ Bessel zeros, to $5.1\times10^{-12}$ for
$R_1$ (whose slowly convergent sum carries an analytic tail
correction) and $\lesssim3\times10^{-15}$ otherwise. At
finite steps the same analytic structure governs the windows: the
poles of the rescaled divisor $f(\rho y)/f(y)$ sit at the $x_n$
for generic step ratio $\rho=T^2/s^2$ (at the ratios
$\rho=(j_{0,k}/j_{0,1})^2$ the numerator vanishes at the dominant
pole $x_1$---precisely the band edges below; cancellations at a
subdominant $x_n$ occur at $(j_{0,k}/j_{0,n})^2$ and shift only
subleading terms),
the residue at $x_1$ carries the factor
$J_0(j_{0,1}\sqrt\rho)$, and for large depth $m$
\[
\operatorname{sign} h_m\to(-1)^m\operatorname{sign}
J_0\bigl(j_{0,1}\sqrt\rho\bigr),
\qquad
h_{m+1}/h_m\to-4/j_{0,1}^2=-0.691660\ldots
\]
(computed $-0.6917$ by $m\approx10$ at $\rho=2,6$; slower near the
band edge at $\rho=14$). The negative windows of
$h_m(\rho)$ converge to \emph{Fourier--Bessel bands} with edges at
$(j_{0,n}/j_{0,1})^2=5.269,\,12.949,\,24.042,\ldots$: the even-$m$
first-window edge marches $3.000\to4.778\to5.261\to5.2689$
($m=2,4,6,8$) against $5.269$, and the odd-$m$ windows converge
onto $(5.269,\,12.949)$, reaching $(5.269,\,12.947)$ at $m=11$.
(The positivity of $h_3$ at every step ratio, which a $2{\times}2$ intuition
would not predict, is exact:
$36\,h_3(\rho)=(\rho-1)(\rho^2-8\rho+19)$ (in the rescaled convention;
the dimensional coefficient carries a further positive $s^6$), and the quadratic factor has
negative discriminant; the first odd-$m$ window opens at $m=5$,
inside the band.) These band statements are dominant-pole
asymptotics checked numerically; they sit below the strength of
Theorems~\ref{thm:k2} and~\ref{thm:rayleigh} and are to be read at
that strength. The sign map of a deep \emph{Jordan} cascade's
divisor thus depends on the depth and the step ratio and on
nothing else in the physics.

%=====================================================================
\section{The deformation dictionary}
\label{sec:dictionary}

The two layers assemble into a dictionary
(Table~\ref{tab:dictionary}).

\begin{table}[t]
\centering\small
\begin{tabular}{@{}>{\raggedright\arraybackslash}p{0.27\textwidth}>{\raggedright\arraybackslash}p{0.08\textwidth}>{\raggedright\arraybackslash}p{0.15\textwidth}>{\raggedright\arraybackslash}p{0.40\textwidth}@{}}
\toprule
deformation & kind & phase & dynamics / divisor image\\
\midrule
Einstein term $m_2^2$ (Stelle~\cite{Stelle1977}) & family &
$\varphi=0$ &
oscillatory; cost $2/\delta$; negativity returns at $\pi/|\delta|$\\
loop width (Merlin~\cite{DonoghueMenezes}) & family &
$\arctan(\gamma/m_2^2)$ & growth $\gamma/2\omega$ per pair
(time-dilated; see text)\\
self-interaction at $k=0$~\cite{IK2014} & corner & $\varphi=\pi$ &
blowup, rate $(\sqrt3/2)\gamma_b$; size-4 block: negativity by
Thm~\ref{thm:k2}\\
non-symplectic lattice / machine noise & generic & random & growing (w.p.\ 1 for unstructured phase-only noise; fully generic $\to1$ (rate $1-O(\sqrt\epsilon)$, bounded phase density); a structured lattice can split real), rate
$\sqrt\epsilon$\\
stacked sectors~\cite{Smilga2014} & --- & --- & mixed mechanisms
(Theorem~\ref{thm:k2})\\
\bottomrule
\end{tabular}
\caption{Deformations of the conformal point, classified by
unfolding kind and phase, with the process-level image. The
oscillatory set is the punctured real axis ($\Delta\neq0$) of the family directions (while
$\omega_2^2$ stays positive) and, the ray $\varphi=0$ (the
single-entry-model reference direction, resolved to within $O(\epsilon)$) of the
generic ones---measure zero in the complex plane either way.}
\label{tab:dictionary}
\end{table}

Two rows need comment. The \emph{Merlin} row records a result that
corrects a natural first guess: one might expect the ghost's
loop-induced width $\gamma$ to impose a detuning floor
$|\omega_1-\omega_2|\gtrsim\Gamma_M$ (the pole width expressed as a
frequency, $\Gamma_M=\gamma/m_r$) and so to detune the theory
off the exceptional line in the deep UV. It does not. The
Donoghue--Menezes pole shift is Lorentz-invariant in $q^2$
($M^2=m_r^2+\ii\gamma$)~\cite{DonoghueMenezes}, so at spatial
momentum $k$ the frequency-space width $\gamma/2E_k$ time-dilates
at exactly the $1/k$ rate of the real detuning $m_2^2/2k$: both
components of the complex detuning vanish together, the condition
number still diverges (suppressed only by the $k$-independent
factor $[1+(\gamma/m_2^2)^2]^{-1/2}\approx1-2\times10^{-4}$ at
$\xi=1$, $N_{\rm eff}\approx20$ in the notation of the revised
\cite{choosing}), and the only change is a fixed
$\approx1^\circ$ tilt of the approach ray into the complex plane.
The dressed near-EP is thus of resonance type---one real and one
complex frequency at every finite $k$---and the saturation a
fixed-linewidth intuition would predict
($\mathrm{cond}_{\max}\sim640\pi/\xi^2N_{\rm eff}$) is an artifact
of freezing the width in frequency. This paragraph's content
appears in the revised companion~\cite{choosing}; we repeat it
here because it is the dictionary's only \emph{tilted} family row.

The \emph{stacked sectors} row is Theorem~\ref{thm:k2} speaking,
with its premise supplied by the companion: the $M$-quantum sector
of the equal-frequency Fock tower carries a single irreducible
Jordan block of size $M{+}1$ (the sector-tower proposition
of~\cite{choosing}; the constructive form of~\cite{Smilga2014}),
so---the sector similarity being a fixed, $t$-independent positive-diagonal one
at every $M$~\cite{choosing}, under which entrywise signs are invariant---the
tower cannot present a pure
mass-creation process; negativity at depth two is a theorem at every sector
$M\ge2$ (for every
sufficiently short step), with generator entry a fixed positive multiple of
$-t^3$ (exactly $-t^3$ in the Jordan basis). Purity is a property of the
minimal cell, and of nothing above it.

%=====================================================================
\section{The photonic readout}
\label{sec:photonic}

Both layers would be measurable with classical light. The
program's experimental proposal~\cite{measurable} extracts
divisibility observables from the intensity transfer matrices
$\hat\Gamma(z)$ of lossy waveguide couplers on the Klauck
platform~\cite{Klauck2019}, and fixes the measurement chain and
error model (fabrication disorder, per-launch calibration, shot
noise, dark floor) that we adopt unchanged here; of the two devices below,
D5 (the unidirectional cascade cell) is new, while D4 adds a new guide-basis
negativity readout to the proposal's existing EP coupler; their error budget
is stated in full in this section and in the public script.

\emph{D4: the dimer at its own EP} ($\gamma=2\kappa$). A reciprocal
coupler's Jordan vector at the EP is the superposition
$(1,-\ii)/\sqrt2$, not a guide, so the guide-basis process is
\emph{not} a cascade and fails by negativity, with onset
${\to}\,1/(\kappa\sqrt2)$ (a convention-free negativity onset, not
the mass-creation $t_c$ of the conventions paragraph), the
continuous EP limit of the
broken-phase onset formula (at $\kappa=0.30$\,cm$^{-1}$: simulated
onset $2.15$\,cm at the $0.25$\,cm divisor step, against the
formula's $2.36$\,cm---the ${\sim}0.2$\,cm shortfall tracking the finite
divisor-step discretization, of which the half-step $\tau/2$ anchor used for
D5 accounts for about half).

\emph{D5: the cascade cell in glass.} Realizing the PU cell in the
beable (guide) basis requires \emph{unidirectional} coupling:
coherent coupling $J$ at synthetic phase $\pi/2$ plus dissipative
coupling $\Gamma_d$ through a common lossy bus gives
$H_{12}=\ii(J-\Gamma_d)$, $H_{21}=-\ii(J+\Gamma_d)$, one-way at
$J=\Gamma_d$ (a Hatano--Nelson asymmetry~\cite{HatanoNelson1996}).
The device shows the mechanism fingerprint on two-length divisor
data: the det-based onset estimator of the original proposal is
\emph{provably blind} (a triangular divisor's determinant is a
product of positive diagonal ratios) while the column sum crosses
one at $z_c=0.91$\,cm ($J=\Gamma_d=0.25$\,cm$^{-1}$, divisor step
$0.25$\,cm; the onset is protocol-dependent, $0.486/2J=0.97$\,cm
in the zero-step limit, and convention-relative in the sense of
Sec.~\ref{sec:divisor}: it is the $\sigma_{\max}$-normalized
observable). The normalization is itself measurable here: for the
triangular device $U^\dagger U$ has diagonal equal to the measured
column sums of $\hat\Gamma$ and off-diagonal modulus
$\smash{\sqrt{\hat\Gamma_{21}\hat\Gamma_{22}}}$, so
$\sigma_{\max}^2(z)$ is the largest eigenvalue of a matrix
assembled from the measured intensities alone---no phase-resolved
characterization is needed. This reconstruction is exact at
triangularity (relative error $6\times10^{-16}$) and, under
asymmetry imperfection, accurate to the first-order data-driven Weyl estimate
$\smash{\sqrt{\hat\Gamma_{11}\hat\Gamma_{12}}}$ set by the measured
residual entry itself (relative $5\times10^{-3}$ at $\eta=3\%$ of
$J$; script block~[1c]). Det-silent but column-sum-loud identifies
the cascade device of this pair; det-crossing with negativity, the
dimer. Detuned
variants ($\delta=0.5$ and $0.8$\,cm$^{-1}$) read the direction:
negativity first returns near
$z^*=\pi/\delta$ (simulated: $6.16$ and $3.81$\,cm at the
$0.25$\,cm divisor step, against $\pi/\delta=6.28$ and
$3.93$\,cm---for $z^*$ the offset is exactly half the pair
spacing: $\pi/\delta-\tau/2=6.158$ and $3.802$\,cm).
Simulated at $10^4$ detected photons per launch
with the proposal's error model ($400$ realizations), using
closed-form, calibration-free estimators for the balanced device
($\Gamma_d$ from the log-slope of the dark-subtracted
$\hat\Gamma_{11}$ ladder; $J+\Gamma_d$ from
$\smash{\sqrt{\hat\Gamma_{21}/\hat\Gamma_{11}}}/z$, in which the
per-launch calibration cancels exactly): $J$ recovered to
$+0.5\%\pm1.7\%$ with recovery correlation $0.74$, printed as a
control---an earlier least-squares model fit returned only its
initialization on the balanced device, so its apparent precisions
were echoes of the injected disorder priors, and the script now
prints that non-recovery ($\mathrm{corr}(\hat\eta,\eta)=0.04$) as
a permanent control; $\hat z_c$ to $\pm0.040$\,cm at the fixed
$0.25$\,cm protocol (the protocol convention itself shifts the
onset by $0.06$\,cm, as stated above, and is part of the quoted
observable); and $z^*$ to $\pm0.22$\,cm about the finite-step
offset (simulated center $-0.10$\,cm vs the $-\tau/2=-0.125$\,cm
anchor; the detuned-device fit that delivers it is itself
controlled, $\mathrm{corr}(\hat\delta,\delta)=0.86$, printed). The asymmetry
imperfection $\eta$ (convention: back-coupling $H_{12}=\ii\eta$ at
fixed $H_{21}=-\ii(J+\Gamma_d)=-0.5\ii$\,cm$^{-1}$) is \emph{not}
resolved by the intensity fit at the $3\%$ design level---its
$\hat\Gamma_{12}$ signal peaks at the dark floor at these photon
numbers---so it is certified rather than fitted, and the certification is
necessarily \emph{two-channel}, because the negativity channel
reads only one sign of $\eta$:
negativity-cleanliness of a $6$\,cm window bounds the real-split
side $\eta/J$ ($\eta>0$) below
the ${\approx}13\%$ boundary (computed: $z^*=6.0$\,cm at $13\%$,
first
in-window negativity at $14\%$ on the $1\%$ grid near the
boundary, boundary
between $13$ and $14\%$; the $z^*$ estimator scans pair base
points at $0.01$\,cm with $0.25$\,cm spacing, the window-negativity
scan at $0.1$\,cm---window-aligned conventions
shift these by up to one grid step), and larger positive
imperfections are
measured directly by the re-entry length (computed: $z^*=6.9$\,cm
at $\eta/J=+10\%$, against $12.7$\,cm---off-chip---at $+3\%$; the
imperfect $H_{12}\neq0$ device is non-triangular, so the triangular
bond rule $z^*=\pi/|\delta|$ no longer governs and onset falls below
it). The
imaginary-split side ($\eta<0$) is negativity-\emph{silent} at any
magnitude (the imaginary split keeps the relevant two-time
determinant $\ge1$ throughout, so the sign flip never triggers;
corroborated computationally, no in-window negativity at $\eta/J=-14\%$
through $-100\%$), so window cleanliness alone cannot bound it;
that side is bounded instead by the Weyl residual of the intensity
reconstruction: the error in $\sigma_{\max}^{2}$ obeys
$\lesssim\smash{\sqrt{\hat\Gamma_{11}\hat\Gamma_{12}}}$
(script block~[1c]; first order in the residual, the $O(\hat\Gamma_{12})$
correction absorbed by the margin below), which is sign-blind in $\eta$---it reads the
measured back-coupling intensity $\hat\Gamma_{12}$ directly---and
bounds the triangularity systematic from the same data. The
proposal's $3\%$
per-launch calibration level is thus met with a factor-four margin (heuristically, to first order in the residual)
against the $13\%$ boundary on the real-split side, and by the
sign-blind intensity residual on the silent side. The
imperfection's \emph{sign} selects the
unfolding direction ($\eta>0$: real split, negativity re-enters at
finite $z$; $\eta<0$: purely imaginary split, clean past
$30$\,cm): even the device's error budget is direction physics.

For feasibility we note that at-EP process physics is
experimentally live (an observation of quantum exceptional points
of an extended Liouvillian on a superconducting circuit appeared
in 2025~\cite{QubitEP2025}---non-Markovian EPs, fittingly for this
program), that Klauck-type arrays reach the required loss rates,
and that the required asymmetric (Hatano--Nelson) coupling has
been realized in photonics: Weidemann et
al.~\cite{Weidemann2020} implemented tunable anisotropic coupling
in a non-Hermitian photonic mesh lattice---the non-Hermitian skin
effect experiment---in the same experimental school as the Klauck
platform. Transferring that coupling asymmetry from the temporal
mesh to a laser-written waveguide pair via a common lossy bus is a
design exercise the proposal's error budget estimates---a single
coupling transfer between demonstrated platforms, itself not yet tested
end-to-end, rather than a new capability.

%=====================================================================
\section{Relation to prior work}
\label{sec:prior}

The Puiseux mathematics of EP splitting is classical~\cite{Kato}
and its physics literature is active---a direction-dependent
Puiseux classification via Hessenberg forms of the perturbation,
with sensing applications~\cite{Mandal2026} (the closest prior
art to Sec.~\ref{sec:map}'s classification, carrying no
divisibility or process content), formalism at exceptional
points~\cite{ChenJu2026}, algebraic characterizations of
non-Hermitian degeneracies~\cite{SayyadStarkov2026}, and, closest
in phrase, ``unitary paths of unfolding'' through a third-order EP
in a solvable model~\cite{BGZ2026}, which exhibits particular
reality-preserving corridors where we classify the
finite-frequency family's deformation directions and attach a
stochastic-mechanism image. The closest
EP--divisibility contact in the literature is the non-Markovian
exceptional-point line~\cite{QubitEP2025,NMEPtheory,WongZengLi2026},
where EPs
emerge in open-system spectra tied to divisibility transitions
(most directly in \cite{WongZengLi2026}, whose channel-level EP
explicitly violates CP-divisibility);
that line classifies neither unfolding directions nor divisor
failure mechanisms. On the gravity
side, Kumar and Marto~\cite{KumarMarto2026} deliberately trade the
Stelle ghost for an inverted-harmonic-oscillator instability and
quantize on that side; in the present language theirs is a
committed choice of the growing side of the same boundary, and the
two sides are related by the unfolding direction of one Jordan
block rather than by separate constructions. In the divisibility
literature, the geometric characterization of stochastic divisors
by Pimenta~\cite{Pimenta2025} is trace-preserving throughout---unit
column sums are structural there---so the column-sum-crossing
mechanism central to this paper and to~\cite{choosing} lies
outside that framework by construction; our sub-stochastic setting
is precisely what makes conditional mass creation available. The
resonant-PU algebraic thread of Fring, Marquette and
Taira~\cite{FMT2026} is engaged in~\cite{choosing}; the present
paper adds the divisor-level image their algebra does not address.
Ilhan and Kovner's two papers~\cite{IK2014,IK2013} supply the
corner instability and the Jordan-block statement respectively;
our use of them is scoped in Sec.~\ref{sec:map}. To our knowledge
no prior work attaches a stochastic divisibility-mechanism image
(a divisor sign-and-column-sum readout) to the directional
unfolding of an exceptional point---the directional Puiseux
classification itself has independent precedent~\cite{Mandal2026},
and the divisor image is what is new here---and no photonic
proposal measures
divisor column sums; we date these claims to the program's public
notebook in the program's hub repository.

%=====================================================================
\section{Discussion}
\label{sec:discussion}

The phase-boundary reading changes what kind of object conformal
gravity is. (We use ``phase boundary'' in the exceptional-point
sense---dynamical regimes separated by a spectral
degeneracy---not in the Landau sense of an order parameter; the
analogy is structural, and we flag it as such.) On this map the pure theory occupies a measure-zero
set: every physical deformation, every regulator, and every
numerical representation sits on the growing, oscillatory, or (for
net-damping representations) decaying side, and the
side is decided by a phase. The theory's notorious fragility at the
equal-frequency point---polynomial growth flipping to exponential
under machine-epsilon perturbations, recorded
in~\cite{choosing}---is not a numerical nuisance but the boundary
itself, and the classification converts anecdotes into a dictionary
(Table~\ref{tab:dictionary}) with a per-row observable.

The divisor theorems give the boundary a laboratory-facing
signature and settle an internal question. Mass-creation-pure
indivisibility---the mechanism that the conformal-gravity paper found
at the conformal cell---is now known to be a property of the minimal Jordan cell
exclusively: depth-three \emph{Jordan} cascades are
negativity-bearing as a matter of Bessel-zero arithmetic, with
exactly alternating generator signs (a reducible three-level
cascade, e.g.\ a $k=2$ cell plus a spectator level, can of course
remain pure---the theorem is about the irreducible block). The universality is striking: the sign map of a
deep Jordan cascade's divisor knows the depth and the step ratio,
and
nothing else about the physics. We do not currently know a
physical principle that singles out the Rayleigh sums here beyond
the calculation itself; the appearance of Fourier--Bessel band
edges $(j_{0,n}/j_{0,1})^2$ in a divisibility question seems to us
genuinely new, and possibly a shadow of something better.

In the program's arc, this paper is the second of the pair the
threads audit projected: the conformal-gravity paper proved what fails \emph{at} the
conformal point; this paper proves what the neighborhood looks
like and how to measure which neighbor nature chose---if a
fourth-order sector is physically realized at all (the depth-two negativity is a
theorem for the formal Jordan block regardless), which remains the program's
standing conditional. The natural continuations are the
finite-time and dressed-propagator questions (the sheet
classification of ghosts versus unstable particles admits a
process-level image we leave to later work), and the boundary
conditions of the de Sitter completion, where the same Jordan
degeneration prices Maldacena's
ghost-removal~\cite{Maldacena2011}; both are planned
as later installments of the program.

\paragraph{Code availability.} All numerical claims are
reproducible from the public scripts in the program
repository---except the Merlin-row loop numbers of
Sec.~\ref{sec:dictionary}, which are analytic computations whose
arithmetic is stated in the revised companion~\cite{choosing}:
\texttt{scripts/p6\_unfolding\_sweep.py} (Sec.~\ref{sec:map}),
\texttt{scripts/p6\_divisor\_gate.py}, \texttt{scripts/p6\_divisor\_theorem.py},
\texttt{scripts/p6\_divisor\_dxd.py}, \texttt{scripts/p6\_cm\_signs.py},
\texttt{scripts/p6\_m2\_channel.py}
(Sec.~\ref{sec:divisor}), and \texttt{scripts/p6\_d4\_device.py}
(Sec.~\ref{sec:photonic}), each printing its controls inline.

\begin{thebibliography}{99}\small

\bibitem{indivisible}
J.~Torkelson,
\emph{Indivisible stochastic processes for non-Hermitian and
higher-derivative quantum systems: detailed balance, broken $\PT$
symmetry, and the explosion of ghosts},
repository \url{https://github.com/embeddetech/p01-indivisible-ghosts}
(2026).

\bibitem{instantaneous}
J.~Torkelson,
\emph{Ghost vacuum decay as a stochastic process: instantaneous
states, the boost-orbit measure, and the foliation that ghosts make
visible}, repository
\url{https://github.com/embeddetech/p03-instantaneous-ghosts} (2026).

\bibitem{conditioned}
J.~Torkelson,
\emph{The price of a post-selected signal: a sharp
survival--distinguishability trade-off, saturated by broken-$\PT$
dynamics}, repository
\url{https://github.com/embeddetech/p04-conditioned-ghosts} (2026).

\bibitem{choosing}
J.~Torkelson,
\emph{Conformal gravity must choose: exceptional-point kinematics,
vacuum decay, and the sample space of Weyl gravity}, repository
\url{https://github.com/embeddetech/p05-choosing-ghosts} (2026).

\bibitem{measurable}
J.~Torkelson,
\emph{Experimental proposal: divisibility onset, generator floor,
and the backflow--divisibility dissociation in a passive $\PT$
photonic coupler}, paper~2 of this program, repository
\url{https://github.com/embeddetech/p02-measurable-ghosts} (2026).

\bibitem{IK2014}
I.~B.~Ilhan and A.~Kovner,
\emph{Conformal gravity redux: Ghost-turned tachyon},
Phys.\ Rev.\ D \textbf{89}, 104015 (2014), arXiv:1403.4664.

\bibitem{IK2013}
I.~B.~Ilhan and A.~Kovner,
\emph{Some comments on ghosts and unitarity: The Pais--Uhlenbeck
oscillator revisited},
Phys.\ Rev.\ D \textbf{88}, 044045 (2013), arXiv:1301.4879.

\bibitem{Kato}
T.~Kato,
\emph{Perturbation Theory for Linear Operators},
2nd ed., Springer (1976).

\bibitem{KreinTheory}
M.~G.~Krein,
\emph{A generalization of certain investigations of A.~M.~Liapunov
on linear differential equations with periodic coefficients},
Dokl.\ Akad.\ Nauk SSSR \textbf{73}, 445 (1950);
R.~S.~MacKay,
\emph{Stability of equilibria of Hamiltonian systems},
in \emph{Nonlinear Phenomena and Chaos}, ed.\ S.~Sarkar,
Adam Hilger, Bristol (1986), pp.~254--270.

\bibitem{Watson}
G.~N.~Watson,
\emph{A Treatise on the Theory of Bessel Functions},
2nd ed., Cambridge University Press (1944).

\bibitem{Smilga2014}
A.~V.~Smilga,
\emph{Exceptional points of infinite order giving a continuous
spectrum},
Int.\ J.\ Theor.\ Phys.\ \textbf{54}, 3900 (2015), arXiv:1409.8450.

\bibitem{DonoghueMenezes}
J.~F.~Donoghue and G.~Menezes,
\emph{Unitarity, stability, and loops of unstable ghosts},
Phys.\ Rev.\ D \textbf{100}, 105006 (2019), arXiv:1908.02416;
\emph{The arrow of causality and quantum gravity},
Phys.\ Rev.\ Lett.\ \textbf{123}, 171601 (2019), arXiv:1908.04170
(the ``Merlin'' coinage).

\bibitem{Barandes}
J.~A.~Barandes,
\emph{The stochastic-quantum correspondence},
Philosophy of Physics \textbf{3}(1), 8 (2025),
doi:10.31389/pop.186,
arXiv:2302.10778;
\emph{The stochastic-quantum theorem}, arXiv:2309.03085.

\bibitem{Stelle1977}
K.~S.~Stelle,
\emph{Renormalization of higher-derivative quantum gravity},
Phys.\ Rev.\ D \textbf{16}, 953 (1977).

\bibitem{Maldacena2011}
J.~Maldacena,
\emph{Einstein gravity from conformal gravity},
arXiv:1105.5632 (2011).

\bibitem{NMEPtheory}
J.-D.~Lin, P.-C.~Kuo, N.~Lambert, A.~Miranowicz, F.~Nori, and
Y.-N.~Chen,
\emph{Non-Markovian quantum exceptional points},
Nat.\ Commun.\ \textbf{16}, 1289 (2025), arXiv:2406.18362.

\bibitem{Klauck2019}
F.~Klauck \emph{et al.},
\emph{Observation of $\PT$-symmetric quantum interference},
Nat.\ Photonics \textbf{13}, 883 (2019), arXiv:1904.08135.

\bibitem{HatanoNelson1996}
N.~Hatano and D.~R.~Nelson,
\emph{Localization transitions in non-Hermitian quantum mechanics},
Phys.\ Rev.\ Lett.\ \textbf{77}, 570 (1996),
arXiv:cond-mat/9603165.

\bibitem{Weidemann2020}
S.~Weidemann, M.~Kremer, T.~Helbig, T.~Hofmann, A.~Stegmaier,
M.~Greiter, R.~Thomale, and A.~Szameit,
\emph{Topological funneling of light},
Science \textbf{368}, 311 (2020).

\bibitem{QubitEP2025}
H.-L.~Zhang, P.-R.~Han, F.~Wu, W.~Ning, Z.-B.~Yang, and
S.-B.~Zheng,
\emph{Experimental observation of non-Markovian quantum exceptional
points},
Phys.\ Rev.\ Lett.\ \textbf{135}, 230203 (2025), arXiv:2503.06977.

\bibitem{ChenJu2026}
W.-M.~Chen and C.-Y.~Ju,
\emph{Non-Hermitian Rayleigh--Schr\"odinger-like perturbation
theory at exceptional point}, arXiv:2606.31279 (2026).

\bibitem{SayyadStarkov2026}
S.~Sayyad and G.~A.~Starkov,
\emph{Characterizing all non-Hermitian degeneracies using algebraic
approaches: Defectiveness and asymptotic behavior},
arXiv:2604.16140 (2026).

\bibitem{BGZ2026}
B.~Bagchi, A.~Ghosh, and M.~Znojil,
\emph{Triple exceptional point with unitary paths of unfolding in a
three-site fermionic Swanson-like model},
Ann.\ Phys.\ (N.Y.) \textbf{492}, 170569 (2026), arXiv:2606.03789.

\bibitem{Mandal2026}
I.~Mandal,
\emph{Puiseux series about exceptional singularities dictated by
symmetry-allowed Hessenberg forms of perturbation matrices},
arXiv:2603.25603 (2026).

\bibitem{KumarMarto2026}
K.~Sravan Kumar and J.~Marto,
\emph{Quantum (quadratic) gravity: replacing the massive tensor
ghost with an inverted harmonic oscillator-like instability},
arXiv:2603.07150 (2026).

\bibitem{Pimenta2025}
L.~S.~Pimenta,
\emph{Divisible and indivisible Stochastic-Quantum dynamics},
arXiv:2505.08785 (2025).

\bibitem{WongZengLi2026}
W.~C.~Wong, B.~Zeng, and J.~Li,
\emph{Non-Markovian exceptional points by interpolating quantum
channels},
npj Quantum Inf.\ \textbf{12}, 63 (2026), arXiv:2507.16049.

\bibitem{FMT2026}
A.~Fring, I.~Marquette, and T.~Taira,
\emph{Spectrum-generating algebra and intertwiners of the resonant
Pais--Uhlenbeck oscillator}, arXiv:2601.20752 (2026).

\end{thebibliography}

\end{document}
